%! Polar Function Graphs — Unit 7, Lesson 7.7 — Group Activity
\documentclass[10pt]{article}
\usepackage{precalculus-article}
\usepackage{precalculus-boxes}
\newcommand{\TallMath}[1]{$\displaystyle #1\rule[-1.4em]{0pt}{3.2em}$}

\begin{document}

\pageheader{Unit 7, Lesson 7.7}{Group Activity}
\namepartnerperiod

\vspace{0.08in}

%----------------------------------------------------------------------
% Context
%----------------------------------------------------------------------
\begin{scenariobox}[Context: Reading Polar Functions]{plumlight}
In polar coordinates, the function $r = f(\theta)$ assigns a radius to each angle.
As $\theta$ increases, the point $(r, \theta)$ traces a curve in the polar plane.
In this activity your group will read polar functions from tables and graphs,
identify key features, and describe what portions of each curve are traced for
given domains. Choose \textbf{one tier} based on your readiness. Everyone should
be prepared to explain their reasoning.
\end{scenariobox}

\vspace{0.08in}

%----------------------------------------------------------------------
% Tier R
%----------------------------------------------------------------------
\begin{tcolorbox}[colback=white, colframe=black!40,
    title=\textbf{Tier R --- Reinforce},
    fonttitle=\bfseries, arc=2mm, left=3mm, right=3mm, top=2mm, bottom=2mm]

\textbf{Directions:} The complete input-output table for $r = 1 + \cos\theta$ is
provided below. Use it to answer the questions.

\vspace{4pt}
\renewcommand{\arraystretch}{1.55}
\begin{tabular}{|c|c|c|c|c|c|c|c|c|c|c|c|c|}
\hline
$\theta$ & $0$ & $\frac{\pi}{6}$ & $\frac{\pi}{3}$ & $\frac{\pi}{2}$ &
           $\frac{2\pi}{3}$ & $\frac{5\pi}{6}$ & $\pi$ &
           $\frac{7\pi}{6}$ & $\frac{4\pi}{3}$ & $\frac{3\pi}{2}$ &
           $\frac{5\pi}{3}$ & $2\pi$ \\
\hline
$r$ & $2$ & $1.87$ & $1.5$ & $1$ & $0.5$ & $0.13$ & $0$ &
      $0.13$ & $0.5$ & $1$ & $1.5$ & $2$ \\
\hline
\end{tabular}

\vspace{8pt}
\textbf{1.}\quad What is the maximum value of $r$? \underline{\hspace{1.5cm}} \quad
At what angle(s) does this occur? \underline{\hspace{2.5cm}}

\vspace{6pt}
\textbf{2.}\quad At what angle is $r = 0$? \underline{\hspace{2cm}} \quad
What does this tell you about the polar curve at that angle? \writeline

\vspace{6pt}
\textbf{3.}\quad On the interval $\theta \in \left[0, \dfrac{\pi}{2}\right]$,
is $r$ increasing or decreasing? \underline{\hspace{2.5cm}}
How do you know from the table? \writeline

\vspace{6pt}
\textbf{4.}\quad On the interval $\theta \in \left[\dfrac{\pi}{2}, \pi\right]$,
is $r$ increasing or decreasing? \underline{\hspace{2.5cm}}

\vspace{6pt}
\textbf{5.}\quad The curve $r = 1 + \cos\theta$ is a limaçon called a
\textbf{cardioid} (heart-shaped). The name comes from the shape. Based on the
table, is the curve symmetric about the $x$-axis? Explain using $r$ values.

\writelines{2}

\end{tcolorbox}

\vspace{0.08in}

%----------------------------------------------------------------------
% Tier A
%----------------------------------------------------------------------
\begin{tcolorbox}[colback=white, colframe=black!40,
    title=\textbf{Tier A --- Apply},
    fonttitle=\bfseries, arc=2mm, left=3mm, right=3mm, top=2mm, bottom=2mm]

\textbf{Directions:} The rectangular graph of $r = \sin(2\theta)$ is shown
below for $\theta \in [0, 2\pi]$. Use it to answer the questions.

\vspace{4pt}
\begin{center}
\begin{tikzpicture}[x=1.05cm, y=1.15cm]
  \draw[linegray, very thin] (-0.2,-1.25) grid[xstep=0.5, ystep=0.5] (6.7,1.25);
  \draw[->] (-0.2,0) -- (6.9,0) node[right]{\small $\theta$};
  \draw[->] (0,-1.3) -- (0,1.35) node[above]{\small $r$};
  \foreach \x/\lbl in {0.785/$\frac{\pi}{4}$, 1.571/$\frac{\pi}{2}$,
                        2.356/$\frac{3\pi}{4}$, 3.142/$\pi$,
                        3.927/$\frac{5\pi}{4}$, 4.712/$\frac{3\pi}{2}$,
                        5.498/$\frac{7\pi}{4}$, 6.283/$2\pi$} {
    \draw (\x,0.06) -- (\x,-0.06);
    \node[below, font=\scriptsize] at (\x,-0.06) {\lbl};
  }
  \foreach \y/\lbl in {-1/{$-1$}, 1/{$1$}} {
    \draw (0.07,\y) -- (-0.07,\y);
    \node[left, font=\scriptsize] at (-0.07,\y) {\lbl};
  }
  \draw[navy, thick, domain=0:6.3, samples=200]
    plot (\x, {sin(2*deg(\x))});
  \node[right, font=\small, navy] at (6.4, 0.4) {$r = \sin(2\theta)$};
\end{tikzpicture}
\end{center}

\vspace{4pt}
\textbf{1.}\quad List all $\theta$ values in $[0, 2\pi]$ where $r = 0$:
\underline{\hspace{6cm}}

\vspace{6pt}
\textbf{2.}\quad On which intervals of $\theta$ is $r > 0$ (positive)?

\underline{\hspace{8cm}}

\vspace{4pt}
On which intervals of $\theta$ is $r < 0$ (negative)?

\underline{\hspace{8cm}}

\vspace{6pt}
\textbf{3.}\quad How many times does $r$ complete a full zero-to-max-to-zero arc
for $\theta \in [0, 2\pi]$? \underline{\hspace{1.5cm}} arcs.

\vspace{4pt}
How many petals does the polar curve $r = \sin(2\theta)$ have?
\underline{\hspace{1.5cm}} petals.

\vspace{6pt}
\textbf{4.}\quad When $r < 0$, the polar point is plotted in the
\emph{opposite} direction from angle $\theta$.
How do the intervals where $r < 0$ contribute to the polar curve?

\writelines{2}

\vspace{6pt}
\textbf{5.}\quad The formula $r = \sin(n\theta)$ has $2n$ petals when $n$ is even.
For $r = \sin(2\theta)$, $n = 2$, so the curve has $2(2) = 4$ petals.
Use the graph to confirm: does the graph show 4 separate zero-to-max-to-zero arcs?
\writeline

\end{tcolorbox}

\vspace{0.08in}

%----------------------------------------------------------------------
% Tier E
%----------------------------------------------------------------------
\begin{tcolorbox}[colback=white, colframe=black!40,
    title=\textbf{Tier E --- Extend},
    fonttitle=\bfseries, arc=2mm, left=3mm, right=3mm, top=2mm, bottom=2mm]

\textbf{Directions:} Two polar curves are shown below. Use the graphs to compare
them and answer the questions.

\vspace{4pt}
\begin{center}
\begin{tikzpicture}[x=1.3cm, y=1.3cm]
  % Polar grid
  \foreach \r in {0.5, 1, 1.5, 2} {
    \draw[linegray, very thin] (0,0) circle (\r);
  }
  \foreach \ang in {0, 30, 60, 90, 120, 150, 180, 210, 240, 270, 300, 330} {
    \draw[linegray, very thin] (0,0) -- ({2.2*cos(\ang)}, {2.2*sin(\ang)});
  }
  \draw[->] (-2.3,0) -- (2.4,0) node[right]{\small $x$};
  \draw[->] (0,-2.3) -- (0,2.4) node[above]{\small $y$};
  % r = 2sin(theta): circle x^2 + (y-1)^2 = 1, radius 1, center (0,1)
  \draw[navy, very thick, domain=0:3.1416, samples=200, smooth]
    plot ({2*sin(deg(\x))*cos(\x r)}, {2*sin(deg(\x))*sin(\x r)});
  % r = 1 + sin(theta): cardioid
  \draw[redacc, very thick, domain=0:6.2832, samples=300, smooth]
    plot ({(1+sin(deg(\x)))*cos(\x r)}, {(1+sin(deg(\x)))*sin(\x r)});
  % Labels
  \node[font=\small, navy] at (0.8, 1.9) {$r = 2\sin\theta$};
  \node[font=\small, redacc] at (1.7, 0.9) {$r = 1+\sin\theta$};
  % Radius labels
  \node[below right, font=\scriptsize] at (1,0) {$1$};
  \node[below right, font=\scriptsize] at (2,0) {$2$};
  \node[above right, font=\scriptsize] at (0,1) {$1$};
  \node[above right, font=\scriptsize] at (0,2) {$2$};
\end{tikzpicture}
\end{center}

\vspace{4pt}
\textbf{1.}\quad For $r = 2\sin\theta$:
\begin{itemize}[leftmargin=1.4em, itemsep=2pt]
  \item Maximum value of $r$: \underline{\hspace{1.5cm}}
  \item $\theta$ value(s) where $r = 0$ (in $[0, 2\pi)$): \underline{\hspace{3cm}}
\end{itemize}

\vspace{4pt}
\textbf{2.}\quad For $r = 1 + \sin\theta$:
\begin{itemize}[leftmargin=1.4em, itemsep=2pt]
  \item Maximum value of $r$: \underline{\hspace{1.5cm}} (occurs at $\theta = $ \underline{\hspace{1.2cm}})
  \item Minimum value of $r$: \underline{\hspace{1.5cm}} (occurs at $\theta = $ \underline{\hspace{1.2cm}})
  \item $\theta$ value(s) where $r = 0$ (in $[0, 2\pi)$): \underline{\hspace{3cm}}
\end{itemize}

\vspace{6pt}
\textbf{3.}\quad Describe in words how the curve $r = 2\sin\theta$ is traced as
$\theta$ goes from $0$ to $2\pi$. When is it moving away from the origin? When
is it moving toward the origin? What happens for $\theta \in (\pi, 2\pi)$?

\writelines{3}

\vspace{6pt}
\textbf{4.}\quad Describe in words how the curve $r = 1 + \sin\theta$ differs
from $r = 2\sin\theta$. Include observations about the maximum radius, the
$\theta$ values where $r = 0$, and the overall shape.

\writelines{3}

\end{tcolorbox}

\end{document}
